% generator_I04.tex -- Prop I.4 (SAS): two triangles equal in every respect.
% Figure geometry derived from Byrne's 1847 construction (public domain).
% Prose authored from the 1847 original, NOT from the CC-BY-SA upstream source.
% Build: lualatex generator_I04.tex  ->  generator_I04.pdf
\documentclass[12pt]{article}
\usepackage{geometry}
\geometry{a5paper, margin=1.5cm}
\usepackage[lining]{ebgaramond}
\usepackage{amssymb}
\usepackage{byrne}
\def\mpPre{u := 1cm; textLabels := false;}

\begin{document}

% FIGURE: two congruent triangles ABC (upper) and DEF (lower).
% ABC: AB red, BC black, CA blue; angles A yellow, B blue, C red.
% DEF: same colors, thinner lines to distinguish.
% DEF is ABC shifted straight down by d.
\defineNewPicture[1/5]{%
    pair A, B, C, D, E, F, d;%
    A := (0, 0);%
    B := A shifted (-5/2u, -7/2u);%
    C := A shifted (1/2u, -3u);%
    d := (0, -4u);%
    D := A shifted d;%
    E := B shifted d;%
    F := C shifted d;%
    % Layer-1 assertions: corresponding sides must be equal (DEF is a translate of ABC)
    if abs(abs(A-B) - abs(D-E)) > 0.01: errmessage "ASSERT: AB != DE"; fi;%
    if abs(abs(C-A) - abs(F-D)) > 0.01: errmessage "ASSERT: CA != FD"; fi;%
    if abs(abs(B-C) - abs(E-F)) > 0.01: errmessage "ASSERT: BC != EF"; fi;%
    % Triangle ABC
    byAngleDefine(B, A, C, byyellow, SOLID_SECTOR);%
    byAngleDefine(A, B, C, byblue,   SOLID_SECTOR);%
    byAngleDefine(B, C, A, byred,    SOLID_SECTOR);%
    draw byNamedAngleResized(BAC, ABC, BCA);%
    byLineDefine(A, B, byred,   SOLID_LINE, REGULAR_WIDTH);%
    byLineDefine(B, C, byblack, SOLID_LINE, REGULAR_WIDTH);%
    byLineDefine(C, A, byblue,  SOLID_LINE, REGULAR_WIDTH);%
    draw byNamedLineSeq(0)(CA,BC,AB);%
    draw byLabelsOnPolygon(B, A, C)(ALL_LABELS, 1);%
    % Triangle DEF (same colors, thin lines)
    byAngleDefine(E, D, F, byyellow, SOLID_SECTOR);%
    byAngleDefine(D, E, F, byblue,   SOLID_SECTOR);%
    byAngleDefine(E, F, D, byred,    SOLID_SECTOR);%
    draw byNamedAngleResized(EDF, DEF, EFD);%
    byLineDefine(D, E, byred,   SOLID_LINE, THIN_WIDTH);%
    byLineDefine(E, F, byblack, SOLID_LINE, THIN_WIDTH);%
    byLineDefine(F, D, byblue,  SOLID_LINE, THIN_WIDTH);%
    draw byNamedLineSeq(0)(FD,EF,DE);%
    draw byLabelsOnPolygon(F, E, D)(ALL_LABELS, 0);%
}

% STATEMENT
\noindent
\begin{minipage}[t]{0.45\linewidth}
\centering
\vspace{0pt}%
\drawCurrentPicture
\end{minipage}%
\hfill
\begin{minipage}[t]{0.52\linewidth}
\vspace{0pt}%
\textbf{Book~I.\enspace Prop.\ IV. Theor.}

\medskip

\noindent\itshape
If two triangles have two sides of the one respectively equal to two sides of
the other, (\drawUnitLine{AB} to \drawUnitLine{DE} and \drawUnitLine{CA} to
\drawUnitLine{FD}) and the angles (\drawAngle{A} and \drawAngle{D}) contained
by those equal sides also equal; then their bases or their sides
(\drawUnitLine{BC} and \drawUnitLine{EF}) are also equal: and the remaining
angles opposite to equal sides are respectively equal
($\drawAngle{B} = \drawAngle{E}$ and $\drawAngle{C} = \drawAngle{F}$):
and the triangles are equal in every respect.

\bigskip

% PROOF
\noindent\upshape
Let two triangles be conceived, to be so placed, that the vertex of one of the
equal angles (\drawAngle{A} or \drawAngle{D}) shall fall upon that of the
other, and \drawUnitLine{AB} to coincide with \drawUnitLine{DE}, then will
\drawUnitLine{CA} coincide with \drawUnitLine{FD} if applied: consequently
\drawUnitLine{BC} will coincide with \drawUnitLine{EF}, or two straight lines
will enclose a space, which is impossible~(ax.~X), therefore
$\drawUnitLine{BC} = \drawUnitLine{EF}$,
$\drawAngle{B} = \drawAngle{E}$ and $\drawAngle{C} = \drawAngle{F}$,
and as the triangles \drawLine[bottom][triangleABC]{AB,CA,BC} and
\drawLine[bottom][triangleDEF]{DE,FD,EF} coincide, when applied,
they are equal in every respect.

\begin{flushright}Q.\ E.\ D.\end{flushright}

\end{minipage}

\end{document}
